\documentclass[12pt,a4paper]{article}
\usepackage{goyabase}

\usepackage{longtable}
\usepackage{calc}

% Pandoc compatibility
\providecommand{\tightlist}{%
  \setlength{\itemsep}{0pt}\setlength{\parskip}{0pt}}
\newcommand{\passthrough}[1]{#1}

\title{\textbf{MongoDB en Docker - UT8 NoSQL}}
\author{Departamento de Informática}
\date{\today}

\usepackage[spanish, es-noshorthands, es-tabla]{babel}
\usepackage[autostyle]{csquotes}
\begin{document}

\maketitle

\section{MongoDB en Docker - UT8 NoSQL}\label{mongodb-en-docker---ut8-nosql}

Este entorno permite trabajar con \textbf{MongoDB}, una base de datos NoSQL \textbf{Documental} diseñada para ofrecer máxima flexibilidad y escalabilidad.

\subsection{El Ecosistema de Datos de GloboTour: El Catálogo Polimórfico}\label{ecosistema}
En GloboTour, usamos MongoDB para gestionar nuestro activo más dinámico: el \textbf{Catálogo de Experiencias}. A diferencia de SQL, aquí cada producto puede tener sus propios atributos sin necesidad de crear cientos de columnas nulas.

\begin{enumerate}
    \item \textbf{Experiencias de Aventura}: 
    \begin{itemize}
        \item \textit{Dato Real:} El \enquote{Salto en Paracaídas} incluye campos únicos como \texttt{altura\_salto: 4000} y \texttt{equipo\_incluido: true}.
    \end{itemize}
    \item \textbf{Experiencias Gastronómicas}: 
    \begin{itemize}
        \item \textit{Dato Real:} La \enquote{Cata de Vinos en La Rioja} incluye un array de \texttt{alergenos: ["sulfitos", "alcohol"]} y el campo \texttt{bodega: "Marqués de Riscal"}.
    \end{itemize}
    \item \textbf{Sistema de Reseñas Incrustadas}:
    \begin{itemize}
        \item \textit{Dato Real:} Las opiniones de los usuarios no están en otra tabla; viven dentro de cada experiencia en un array \texttt{reviews: [{user: "paco", score: 5, comment: "Increíble"}]}.
    \end{itemize}
\end{enumerate}

\subsection{Componentes}\label{componentes}

\begin{enumerate}
\def\labelenumi{\arabic{enumi}.}
\tightlist
\item
  \textbf{MongoDB Server}: Motor de base de datos en el puerto
  \texttt{27017}.
\item
  \textbf{Mongo Express}: Interfaz web en el puerto
  \texttt{8081}.
\end{enumerate}

\hypertarget{instrucciones-de-uso}{%
\subsection{Instrucciones de uso}\label{instrucciones-de-uso}}

\subsubsection{Configuración del Entorno (Docker Compose)}\label{configuracion-del-entorno}

Este archivo define los servicios necesarios para la práctica:

\begin{lstlisting}[language=YAML]
services:
  mongodb:
    image: mongo:latest
    container_name: mongodb-server
    ports:
      - "27017:27017"
    volumes:
      - mongo_data:/data/db
      - ./src:/src
    restart: always

  mongo-express:
    image: mongo-express:latest
    container_name: mongo-express
    ports:
      - "8081:8081"
    environment:
      - ME_CONFIG_MONGODB_SERVER=mongodb
      - ME_CONFIG_BASICAUTH_USERNAME=admin
      - ME_CONFIG_BASICAUTH_PASSWORD=pass
    depends_on:
      - mongodb
    restart: always

volumes:
  mongo_data:
\end{lstlisting}

\subsubsection{Puesta en marcha}\label{puesta-en-marcha}

Para arrancar el servicio:

\begin{lstlisting}[language=bash]
docker compose up -d
\end{lstlisting}

Para importar los datos de ejemplo:

\begin{lstlisting}[language=bash]
docker exec -i mongodb-server mongoimport --db globotour --collection experiencias --file /src/experiencias.json --jsonArray
\end{lstlisting}

Para acceder a la consola (Mongo Shell):

\begin{lstlisting}[language=bash]
docker exec -it mongodb-server mongosh
\end{lstlisting}

Una vez dentro de Mongo Shell, puedes indicar la base de datos a
utilizar

\begin{lstlisting}[language=bash]
test> use globotour
globotour> 
\end{lstlisting}

Para acceder a la interfaz gráfica: - URL:
\texttt{http://localhost:8081} - Usuario:
\texttt{admin} - Contraseña:
\texttt{pass}

\end{document}
